\section{An elliptic quotient of the Fermat quartic}
\label{sec:finite-geometry}

The specialization criterion asks for a divisor-class homomorphism
that kills the omitted points but detects a weighted sum of the
chosen specializations. We construct it from a degree-two quotient
of the Fermat quartic. The square of the cyclotomic character takes
values $\pm1$; the quotient will turn precisely these signs into
elliptic negation.

Set
\[
    k_0=\F_{5^4},\qquad k=\F_{5^8}.
\]
Over $\F_5$, let $D$ be the projective plane curve
\[
    X^4+Y^4=Z^4,
\]
and let $E$ be the smooth projective model of
\[
    w^2=s^4+1.
\]
The partial derivatives show that $D$ is smooth, including after any
extension of the ground field. It is geometrically integral: distinct
projective plane components would intersect and produce a singular point.
The polynomial $s^4+1$ has four simple roots. At infinity the coordinates
$q=1/s$ and $r=w/s^2$ give $r^2=1+q^4$, with two unramified points
$r=\pm1$ above $q=0$. Thus the double cover $E\to\mathbf P^1_s$ has
exactly four branch points; the degree-two Hurwitz formula gives
genus one.

There is a finite morphism of degree two
\begin{equation}
\label{eq:elliptic-quotient}
    h:D\longrightarrow E,
    \qquad s=\frac XY,\qquad w=\frac{Z^2}{Y^2}.
\end{equation}
This is the familiar passage from a diagonal quartic to a genus-one
quotient by a sign involution; related quotient constructions appear
in \cite[\S1, equations~(1.1)--(1.3)]{FlynnWetherell2001}.
We verify the particular map needed here. On $Y\ne0$, write $v=Z/Y$. The extension of function fields is
obtained by adjoining $v$ with $v^2=w$; the nontrivial involution
$v\mapsto-v$ proves that its degree is two. The resulting rational map
extends to the smooth projective curves and is finite. It maps the four
points of $D$ on $Z=0$ to the four branch points of $E\to\mathbf P^1$.

The eighth roots of unity lie in $k_0$, since $8\mid(5^4-1)$.
Thus all four branch points are $k_0$-rational. Choose one, denoted by $o$,
as the origin of $E_{k_0}$. Every fiber of $E\to\mathbf P^1$ is linearly
equivalent to $2(o)$, so the involution $w\mapsto-w$ is negation for this
group law, using the identification of a degree-zero divisor class
with its point on a genus-one curve with chosen origin
\cite[Proposition~III.3.4]{Silverman2009}. In particular, every branch point represents a point of
$E[2](k_0)$.

We now construct the promised homomorphism on divisor classes. Write
$D_k$ and $E_k$ for the base changes to $k$, and continue to denote the
quotient map by $h:D_k\to E_k$. The origin $o\in E(k_0)$ defines the
group law on $E_k$. Put
\[
 Z_\infty=D_k\cap\{Z=0\}.
\]
Define
\begin{equation}\label{eq:elliptic-detector}
 \begin{split}
 \psi:\Pic(D_k)&\longrightarrow E(k)/5E(k),\\
 [D']&\longmapsto
 [h_*D'-\deg(D')\,(o)]\pmod{5E(k)}.
 \end{split}
\end{equation}
Every line bundle on the smooth curve $D_k$ has a nonzero rational
section and hence a representing divisor $D'$. On a closed point $r$,
finite divisor pushforward is
$h_*(r)=[\kappa(r):\kappa(h(r))](h(r))$. Its degree is
$[\kappa(r):\kappa(h(r))][\kappa(h(r)):k]=[\kappa(r):k]$,
so it preserves degree over $k$.
It takes principal divisors to principal divisors by the norm formula
\cite[Tag \texttt{02RT}]{StacksProject}. The identification
$\Pic^0(E_k)\cong E(k)$ sends $[(R)-(o)]$ to $R$
\cite[Proposition~III.3.4 and Remark~III.3.5.1]{Silverman2009}.
Thus \eqref{eq:elliptic-detector} is a well-defined homomorphism on the
entire Picard group, including classes represented by nonrational closed
points. For a $k$-rational point $r$ of $D_k$, it gives
$\psi([r])=h(r)$ modulo $5E(k)$.

Each point of $Z_\infty$ is $k$-rational and maps to a branch point of
$E_k\to\mathbb P^1_k$. Its degree-zero class is two-torsion, and hence
zero modulo five. We have therefore proved
\begin{equation}\label{eq:detector-at-infinity}
 \psi([r])=0\qquad(r\in Z_\infty).
\end{equation}
It remains to find an orbit whose weighted image under $\psi$ is nonzero.

For $a\in\F_5^\times$, define an automorphism of $D$ by
\[
    g_a(X:Y:Z)=(aX:aY:Z).
\]
We will use the following orbit, with coefficients interpreted in the
group $E(k)/5E(k)$.

\begin{lemma}
\label{lem:orbit}
There are points $Q\in E(k_0)$ and $P\in D(k)$ such that
\[
    Q\notin5E(k),\qquad h(P)=Q,\qquad Z(P)\ne0,
\]
and
\begin{equation}
\label{eq:weighted-orbit}
 \begin{aligned}
    \sum_{a\in\F_5^\times}a^2\psi([g_aP])
      &=\sum_{a\in\F_5^\times}a^2h(g_aP)\\
      &=4Q\ne0\quad\text{in }E(k)/5E(k).
 \end{aligned}
\end{equation}
\end{lemma}

\begin{proof}
The underlying curve $E$ has four points over $\F_5$. At $s=0$ there are two affine
points, while at each nonzero $s$ the value $s^4+1=2$ is a nonsquare.
There are also two rational points at infinity, since the leading
coefficient of the quartic is a square. Choosing any rational point as
origin over $\F_5$, the Frobenius polynomial is therefore
$T^2-2T+5$, and its roots compute extension-field point counts
\cite[Theorem~V.2.3.1]{Silverman2009}. If $a_m$ is the sum of the $m$th powers of its roots, then
\[
    a_0=a_1=2,\qquad a_m=2a_{m-1}-5a_{m-2}.
\]
It follows that $a_4=-14$ and
\[
    \#E(k_0)=5^4+1-a_4=640.
\]
This count is independent of the choice of origin. Hence, with origin
$o$, the finite group $E(k_0)$ has a point $Q\notin5E(k_0)$.

Such a point remains nondivisible by $5$ over $k$. In fact, if $Q=5R$
with $R\in E(k)$, the trace for the quadratic extension $k/k_0$, using
the group law defined over $k_0$, would give
\[
    2Q=5\bigl(R+\operatorname{Frob}_{k_0}(R)\bigr)
        \quad\text{in }E(k_0).
\]
Multiplication by $2$ is invertible on $E(k_0)/5E(k_0)$, a contradiction.
Every branch point is $2$-torsion and therefore divisible by $5$, so
$Q$ is not a branch point. A fiber of the degree-two morphism $h$ has a
point over an extension of $k_0$ of degree at most two. Since $k$ is its
quadratic extension, we may choose $P\in D(k)$ over $Q$. The description
of $h$ on $Z=0$ gives $Z(P)\ne0$.

Finally, on the dense chart $Y\ne0$ the morphism $h\circ g_a$
fixes $s$ and multiplies $w$ by $a^{-2}$. It therefore agrees with
$h$ when $a^2=1$ and with $[-1]\circ h$ when $a^2=-1$. These are
identities of morphisms on all of $D$: they extend from the dense
chart because $D$ is integral and the target is separated. This
also covers the possibility that $Q$ is a point at infinity. Thus
\[
    h(g_aP)=
    \begin{cases}
        Q,&a^2=1,\\
        -Q,&a^2=-1.
    \end{cases}
\]
Each weighted summand in \eqref{eq:weighted-orbit} equals $Q$ modulo
$5$. Their sum is $4Q$, which is nonzero because multiplication by $4$
is invertible on $E(k)/5E(k)$.
\end{proof}

We fix $Q,P$ as in the lemma. The homomorphism $\psi$ kills
$Z_\infty$ by \eqref{eq:detector-at-infinity} and detects the weighted
orbit of $P$ by \eqref{eq:weighted-orbit}. The next section realizes
these four specializations as the conjugates of one horizontal divisor.
Curves and morphisms in residue characteristic will henceforth be
viewed over $k$.

\section{The regular local ring and its split divisor}
\label{sec:local-geometry}

We next realize the four-point orbit as the specialization of a split
horizontal divisor. The construction has three parts: choose a
cyclotomic uniformizer with controlled conjugates, map the ramified
regular local ring onto the cyclotomic valuation ring, and split this divisor
by a flat scalar extension.

Let $H_0/\Q_5$ be the unramified extension of degree eight, and let $V$
be its ring of integers, with residue field identified with $k$.
Retain the cyclotomic notation $\zeta$, $u=\zeta-1$,
$\Delta=\Gal(\Q(\zeta)/\Q)$, and
$\chi:\Delta\xrightarrow{\sim}\F_5^\times$.
Choose an embedding of $\Q(\zeta)$ into an algebraic closure of $H_0$,
and put
\[
    L=H_0(\zeta),\qquad \OO=\text{the ring of integers of }L.
\]
The translated cyclotomic equation
\[
    u^4+5u^3+10u^2+10u+5=0
\]
is Eisenstein over $V$. Consequently $L/H_0$ is totally ramified of
degree four, $u$ is a uniformizer of $\OO$, and the residue field is
$k$. Restriction identifies $\Gal(L/H_0)$ with $\Delta$.

There is a uniformizer $t\in\OO$ satisfying
\begin{equation}
\label{eq:cyclotomic-uniformizer}
    t^4=-5,\qquad \frac tu\equiv1\pmod u,\qquad
    \frac{\delta(t)}t\in\mu_4(V),\qquad
    \overline{\frac{\delta(t)}t}=\chi(\delta).
\end{equation}
To construct it, the cyclotomic equation gives
\[
    \frac{u^4}{-5}=1+2u+2u^2+u^3.
\]
The inverse of this unit has a fourth root congruent to $1$ modulo $u$
by Hensel's lemma; multiply $u$ by that root. The quotient
$\delta(t)/t$ is a fourth root of unity and has residue $\chi(\delta)$,
because $\delta(u)/u$ has that residue. All fourth roots of unity lie
in $V$, again by Hensel's lemma.

We also have
\begin{equation}
\label{eq:monogenic-dvr}
    V[t]=\OO.
\end{equation}
For completeness, $T^4+5$ is Eisenstein, so $V[t]$ is finite free over
$V$ and has fraction field $L$. Reduction modulo $5$ shows that it is
local, with maximal ideal $(5,t)=(t)$. Its dimension is one, hence it
is a DVR. It is therefore integrally closed in $L$, and every element
of $\OO$, being integral over $V$, belongs to $V[t]$.

Write $P=(\bar X_0:\bar Y_0:1)$. Hensel's lemma supplies lifts
$X_0,Y_0\in V$ with
\begin{equation}
\label{eq:lifted-point}
    X_0^4+Y_0^4=1.
\end{equation}
Indeed, at least one of $\bar X_0,\bar Y_0$ is nonzero, and the
corresponding partial derivative is a unit.

\begin{lemma}
\label{lem:local-ring}
The ring
\begin{equation}
\label{eq:counterexample-ring}
    A=\bigl(V[x,y]/(5+x^4+y^4)\bigr)_{(5,x,y)}
\end{equation}
is a regular local domain of dimension two. The assignment
\begin{equation}
\label{eq:closed-divisor-map}
    A\longrightarrow\OO,\qquad
    x\longmapsto tX_0,\qquad y\longmapsto tY_0
\end{equation}
is a surjective local homomorphism with nonzero height-one kernel.
\end{lemma}

\begin{proof}
The regular local ring $V[x,y]_{(5,x,y)}$ has dimension three. The
equation $5+x^4+y^4$ is congruent to $5$ modulo the square of its
maximal ideal, so it is part of a regular system of parameters.
Its quotient is regular local of dimension two and is a domain
\cite[Tags \texttt{00NQ} and \texttt{00NP}]{StacksProject}.

Equations \eqref{eq:cyclotomic-uniformizer} and
\eqref{eq:lifted-point} show that the assignment respects the defining
equation. Every denominator inverted in \eqref{eq:counterexample-ring}
has nonzero constant residue at $x=y=0$, hence maps to a unit of
$\OO$. Thus the homomorphism is defined and local. At least one of
$X_0,Y_0$ is a unit of $V$, so its image contains $t$. Surjectivity
follows from \eqref{eq:monogenic-dvr}.

The element $Y_0x-X_0y$ belongs to the kernel and has nonzero image in
$\mathfrak m_A/\mathfrak m_A^2$. In particular the kernel is nonzero.
It is a nonmaximal prime, since the quotient $\OO$ is a DVR, and
therefore has height one in the two-dimensional local domain $A$.
Since $Y_0x-X_0y$ is a regular parameter, it generates a height-one
prime and hence generates this kernel.
\end{proof}

We have obtained the closed divisor that will support the integral
class. To detect its pushforward, we now identify the generic curve
on which this support separates into conjugate points.

Let
\[
    \Dcal=\Proj\bigl(\OO[X,Y,Z]/(X^4+Y^4-Z^4)\bigr),
    \qquad C=\Dcal_L,
    \qquad T=\Spec(A\otimes_V L).
\]
The model $\Dcal$ is projective. Over an algebraic closure of either
$L$ or $k$, the three partial derivatives $4X^3,4Y^3,-4Z^3$ cannot
vanish at a projective point. Thus each geometric fiber is smooth of
dimension one. Each is also integral: smoothness makes it reduced,
while two distinct plane components would intersect and create a
singularity. The same Jacobian calculation on the relative affine charts,
where $4$ is invertible, proves smoothness over $\OO$
\cite[Tag \texttt{00TE}]{StacksProject}. Hence $\Dcal$ is regular.
Flatness over the DVR excludes a component supported in the special
fiber, and the generic fiber is integral; therefore $\Dcal$ is an
integral surface with special fiber $D_k$.

\begin{lemma}
\label{lem:splitting}
The scheme $T$ is regular, integral, Noetherian, and one-dimensional,
and its morphism to $\Spec A$ is flat. The substitutions
\[
    X/Z=x/t,\qquad Y/Z=y/t
\]
identify $T$ with the spectrum of a localization of the coordinate
ring of the chart $Z\ne0$ in $C$.

The base change to $T$ of the closed immersion
$\Spec\OO\to\Spec A$ in \eqref{eq:closed-divisor-map} is the disjoint
union of four $L$-rational closed points $p_\delta$, indexed by
$\delta\in\Delta$. Their coordinates are
\begin{equation}
\label{eq:split-points}
    p_\delta:
    \quad (X/Z,Y/Z)=
    \left(\frac{\delta(t)}tX_0,
          \frac{\delta(t)}tY_0\right).
\end{equation}
Their specializations on $\Dcal$ are $g_{\chi(\delta)}P$.
The factor indexed by $\delta$ is obtained from $\OO$ by the map
$a\mapsto\delta(a)$ into $L$.
\end{lemma}

\begin{proof}
The $V$-module $L$ is flat, since $V\to H_0$ is a localization and
$L$ is a vector space over $H_0$. The substitution $x=tX$, $y=tY$
transforms $5+x^4+y^4=0$ into $X^4+Y^4=1$. Tensoring and localizing
therefore give the asserted description of $T$. Its coordinate ring
is a localization of the coordinate ring of a smooth geometrically
integral affine curve over $L$; it is regular, integral, and
Noetherian, of dimension at most one.

Since $5$ is invertible in the second tensor factor, the coordinate
ring of the base-changed divisor is
\begin{equation}
\label{eq:divisor-tensor-splitting}
    \OO\otimes_V L
      =L\otimes_{H_0}L
      \xrightarrow{\sim}\prod_{\delta\in\Delta}L,
    \qquad a\otimes b\longmapsto(\delta(a)b)_\delta.
\end{equation}
Here $\OO[1/5]=L$, and the last isomorphism is the usual splitting
of a finite Galois extension. The images of $x$ and $y$ on the
$\delta$-factor are $\delta(t)X_0$ and $\delta(t)Y_0$, giving
\eqref{eq:split-points}. Every inverted denominator evaluates to a
unit of $\OO$, so these factors indeed define points of $T$.
They are distinct: the ratios $\delta(t)/t$ are the four different
fourth roots of unity, and at least one of $X_0,Y_0$ is a unit.
Thus $T$ has closed points and has dimension one.

The coordinates in \eqref{eq:split-points} are integral over $\OO$
and satisfy the Fermat equation, so they define sections of $\Dcal$.
Reducing \eqref{eq:cyclotomic-uniformizer} gives their stated
specializations. The last assertion follows directly from the
factor maps in \eqref{eq:divisor-tensor-splitting}.
\end{proof}
